\section{Conclusion}

This paper has presented a philosophical framework that explores the implications of taking experience as fundamental. By proposing a model of reality as an evolving experiential mesh, a network of interwoven and evolving vibes so-to-speak, it attempts to sketch how consciousness, agency, and physical properties might potentially emerge from experiential dynamics.

If this framework were correct, it would suggest interesting reinterpretations of familiar concepts. Consciousness might be understood as universal, with complexity varying by organization. Individual identity could be seen as patterns rather than fixed entities. Death might represent transformation rather than termination. Physical laws could be understood as regularities in experiential dynamics. Concepts from various wisdom traditions might find natural expression within this framework.

To understand the complete picture: vibes are fundamental experiential units linked by mutual awareness, forming the vibe mesh. Everything updates in discrete beats. Coherent vibe groups form minds at every scale, from single vibes to cosmic structures. The boundaries between self/other and mind/matter depend on perspective within the mesh. Intelligence navigates between states. Learning updates relational patterns. Information becomes knowledge when structured, meaning when useful, memory when stable. Reality is this infinite mesh of experiencing nodes, oscillating between pleasure and pain, creating all existence through their interactions.

While these ideas remain speculative, they demonstrate the potential fertility of experience-based ontologies. This framework offers one possible approach to the hard problem of consciousness, suggesting that perhaps the problem arises from starting with the wrong primitive concepts. Future work could explore whether these ideas can be developed into testable hypotheses or whether they might inspire new approaches in consciousness studies, physics, or artificial intelligence research.
